% Generated by Listas/make_statements.py from pratica-lecture-06.md, item 25. Do not edit.
\textbf{Source.} Cabral (2017), \emph{Introduction to Industrial Organization}, 2nd ed., Ch. 9 Challenging exercises, Exercise 9.19 ``Secret price cuts'', p.~245 {[}PDF 257{]}. Reproduced verbatim, with the margin note (s) it cites (the book's ``price cups'' kept). Answer in the source: none in the book.



\begin{quote}

\textbf{9.19. Secret price cuts.} This exercise formalizes the model of secret price cups presented in Section 9.2. Suppose that all consumers are willing to pay \(u\) for the (homogeneous) product sold by two duopolists. In each period, demand can be high (probability \(1 - \alpha\)) or low (probability \(\alpha\)). When demand is high, \(h = 1\) units can be sold at price \(u\) (or any lower price). When demand is low, only \(l = 0 < h\) units can be sold. The probability that demand is high or low in each period is independent of what it was in the previous period. Moreover, \emph{firms are unable to observe the state of market demand}; all they can observe is whether \emph{their own} demand is high or low. Finally, for simplicity, assume that production costs are zero.



Consider the following equilibrium strategies. Firms start by setting \(p = u\). If they receive a positive demand (namely, \(\frac{1}{2}\)), then they continue to set \(p = u\), that is, they remain in the ``cooperative phase''. If however one of the firms (or both) receives zero demand, then both firms enter into a ``price war:'' they set \(p = 0\) during \(T\) periods and, after this period, revert to \(p = u\) again (the cooperative phase).\(^{s}\) Let \(V\) be the net present value of a firm in equilibrium (starting in a period where collusion takes place).



\begin{enumerate}

\def\labelenumi{(\alph{enumi})}

\item

  Show that:

  \[V = (1 - \alpha)\left(\frac{u}{2} + \delta V\right) + \alpha\,\delta^{T+1} V\]

\item

  Show that the best a firm can do by deviating is to get:

  \[V' = (1 - \alpha)\,u + \delta^{T+1} V\]

\item

  Show that the condition that the prescribed strategy constitutes an equilibrium is given by:

  \[1 \le 2\,(1 - \alpha)\,\delta + (2\,\alpha - 1)\,\delta^{T+1}\]

\item

  What is then the optimal equilibrium?

\end{enumerate}



\emph{s.} Notice that, if a firm receives zero demand, then it is common knowledge that a price war is going to start, that is, it is common knowledge that one of the firms receives zero demand. In fact, either demand is low, in which case both firms receive zero demand, or one of the firms deviates from \(p = u\), in which case the deviating firm knows that the rival receives zero demand.

\end{quote}

